\documentclass{article}
\usepackage{amsmath}
\usepackage{enumitem}

\begin{document}
    
    \section{Self-Awareness}
        
        Self-awareness generally refers to a system's capacity to represent, monitor, and reason about its own internal states, processes, and actions. In cognitive science and artificial intelligence, it is typically considered a higher-order cognitive capability emerging from several underlying components and organized into different forms and levels.
    
    
    \section{Components of Self-Awareness}
        
        Self-awareness is composed of several interacting cognitive mechanisms:
        
        \begin{enumerate}[label=\arabic*.]
\item \textbf{Self-representation}

The ability of an agent to maintain an internal model of itself. This includes representations of its body, capabilities, goals, and internal states.

\item \textbf{Self-monitoring}

Continuous observation of internal processes such as perceptions, actions, decisions, and internal states.

\item \textbf{Self-evaluation}

The ability to assess one's own performance, actions, or internal states against goals, expectations, or norms.

\item \textbf{Self-regulation}

The ability to modify behavior or internal processes based on self-evaluation.

\item \textbf{Meta-representation}

Representation of one's own cognitive processes (thinking about thinking).

\item \textbf{Temporal self-continuity}

Awareness that the self persists across time (past, present, and anticipated future states).
        \end{enumerate}
        
        These mechanisms together allow an agent to treat itself as an object of cognition.
    
    
    \section{Kinds of Self-Awareness}
        
        Different types of self-awareness have been identified in psychology, philosophy, and cognitive science.
        
        \subsection{Bodily Self-Awareness}
            
            Awareness of one's own physical body and its position in space.
            
            Examples:
            \begin{itemize}
                \item proprioception
                \item body ownership
                \item sense of agency over movements
            \end{itemize}
        
        \subsection{Perceptual Self-Awareness}
            
            Awareness that one's perceptions belong to oneself.
            
            Example:
            
            ``I am the one seeing this object.''
        
        \subsection{Reflective Self-Awareness}
            
            The ability to reflect on one's own mental states.
            
            Examples:
            
            \begin{itemize}
                \item awareness of beliefs
                \item awareness of emotions
                \item awareness of intentions
            \end{itemize}
        
        \subsection{Narrative Self-Awareness}
            
            Awareness of oneself as a temporally extended entity with a personal history.
            
            This includes autobiographical memory and personal identity.
        
        \subsection{Social Self-Awareness}
            
            Awareness of how one appears to others and how one is evaluated socially.
            
            Examples:
            
            \begin{itemize}
                \item embarrassment
                \item reputation management
                \item theory of mind about how others perceive the self
            \end{itemize}
    
    
    \section{Levels of Self-Awareness}
        
        Self-awareness can also be organized into hierarchical levels of cognitive complexity.
        
        \subsection{Level 0: No Self-Model}
            
            Systems react to stimuli but have no representation of themselves.
            
            Examples:
            
            \begin{itemize}
                \item simple reflex systems
                \item basic reactive robots
            \end{itemize}
        
        \subsection{Level 1: Bodily Awareness}
            
            The agent represents its body and physical state.
            
            Examples:
            
            \begin{itemize}
                \item proprioception
                \item spatial self-localization
            \end{itemize}
        
        \subsection{Level 2: Agent Awareness}
            
            The system represents itself as the cause of actions.
            
            Examples:
            
            \begin{itemize}
                \item sense of agency
                \item action ownership
            \end{itemize}
        
        \subsection{Level 3: Self-Monitoring}
            
            The agent observes its own internal processes.
            
            Examples:
            
            \begin{itemize}
                \item monitoring internal states
                \item detecting errors in one's own actions
            \end{itemize}
        
        \subsection{Level 4: Meta-Cognitive Awareness}
            
            The agent can reason about its own cognition.
            
            Examples:
            
            \begin{itemize}
                \item evaluating confidence in decisions
                \item reasoning about knowledge or ignorance
            \end{itemize}
        
        \subsection{Level 5: Narrative or Reflective Self}
            
            The agent constructs a temporally extended model of itself.
            
            Examples:
            
            \begin{itemize}
                \item autobiographical memory
                \item long-term self-identity
            \end{itemize}
    
    
    \section{Relationship Between the Concepts}
        
        The hierarchy can be summarized as:
        
        \[
            \text{Body Awareness}
            \rightarrow
            \text{Agency Awareness}
            \rightarrow
            \text{Self-Monitoring}
            \rightarrow
            \text{Meta-Cognition}
            \rightarrow
            \text{Reflective/Narrative Self}
        \]
        
        Self-awareness emerges when a system forms internal models of itself and uses these models to monitor, evaluate, and regulate its behavior.

\end{document}